\labsection{9}{Эволюция эвристик с помощью большой языковой модели}

\labpart{Цель работы}
Изучить эволюционный поиск программ, в котором оператором мутации служит большая языковая модель
(подход FunSearch); реализовать безопасный оценщик сгенерированного кода и мини-FunSearch для задачи
онлайн-упаковки в контейнеры; сравнить найденные эвристики с классическими правилами и с параметрической
эвристикой, настроенной CMA-ES, и оценить их обобщение на новые данные.

\labpart{Теоретические сведения}
Подраздел~\ref{sec:llm-evolution} (FunSearch, требования к оценщику, онлайн-упаковка в контейнеры),
подраздел~\ref{sec:cmaes} (CMA-ES для настройки параметров эвристики), подраздел~\ref{sec:generalization}
(обобщение и переобучение) и подраздел~\ref{sec:ga-method} (методика сравнения стохастических алгоритмов).

\labpart{Постановка задачи}
Предметы с целыми размерами, сгенерированными из распределения варианта (таблица~\ref{tab:l9-var}) и
ограниченными отрезком $[1; C]$, поступают по одному. Эвристика задаётся функцией приоритета
\begin{lstlisting}
def priority(item, bins, capacity):   # bins: свободное место в контейнерах, куда item помещается;
    return -(bins - item)             # последний элемент - новый пустой контейнер (Best Fit)
\end{lstlisting}
и предмет кладётся в контейнер с наибольшим приоритетом. Качество эвристики~--- среднее по наборам
превышение нижней границы, \%:
$$
E = \frac{100}{K} \sum_{k=1}^{K} \frac{m_k - L_1^{(k)}}{L_1^{(k)}}, \qquad
L_1^{(k)} = \Bigl\lceil \sum_i s_i^{(k)} / C \Bigr\rceil,
$$
где $m_k$~--- число контейнеров, использованных на $k$-м наборе. Поиск ведётся на обучающих наборах
($K = 5$ наборов по 1000 предметов, seed 0--4), итог проверяется на тестовых ($K = 5$ по 5000
предметов, seed 100--104).

\emph{Часть А.} Реализовать окружение упаковки и классические правила First Fit, Best Fit, Worst Fit.
Сформулировать параметрическую эвристику~--- Best Fit со штрафом $a$ за мелкий бесполезный остаток
меньше $t$ и бонусом $b$ за точное заполнение:
$$
p(r) = -r - a\,[0 < r < t] + b\,[r = 0], \qquad r = \text{bins} - \text{item},
$$
где $[\cdot]$~--- индикатор условия, и подобрать параметры $(a, t, b)$ своей CMA-ES из ЛР~№8.

\emph{Часть Б.} Реализовать мини-FunSearch: база программ с оценками, подсказка с двумя версиями
функции, выбранными турниром и упорядоченными по возрастанию качества, генерация новой версии языковой
моделью, изолированная оценка кода. Начальная программа~--- Best Fit, бюджет~--- 60 вызовов модели
(15 волн по 4 параллельных вызова).

\emph{Часть В.} Проверить обобщение: оценить пять лучших эвристик на тестовых наборах и на варианте
переноса~--- том же распределении с другой ёмкостью (вариант $k + 4$ для $k \le 4$ и $k - 4$ для $k > 4$).
Выполнить абляцию лучшей эвристики: удалять из кода по одному правилу и измерять вклад каждого.

\begin{table}[!htb]
\caption{Варианты заданий к лабораторной работе №9 и контрольные значения $E$ на обучающих наборах для
генератора из приложения~\ref{app:listings}}
\label{tab:l9-var}
\centering\small
\begin{tabular}{clccc}
\toprule
№ & Распределение размеров предметов & $C$ & First Fit, \% & Best Fit, \% \\
\midrule
1 & равномерное $U[20; 100]$                      & 150 & 5,40  & 4,85 \\
2 & нормальное $N(50; 15)$                         & 100 & 5,82  & 4,37 \\
3 & Вейбулла, масштаб 45, форма 3                   & 100 & 5,25  & 4,76 \\
4 & смесь $0,5\,N(30; 5) + 0,5\,N(65; 5)$           & 120 & 10,19 & 9,49 \\
5 & равномерное $U[20; 100]$                       & 120 & 6,73  & 4,94 \\
6 & нормальное $N(50; 15)$                         & 150 & 4,40  & 4,22 \\
7 & Вейбулла, масштаб 45, форма 3                   & 150 & 2,82  & 2,75 \\
8 & смесь $0,5\,N(30; 5) + 0,5\,N(65; 5)$           & 150 & 6,22  & 6,09 \\
\bottomrule
\end{tabular}
\end{table}

\labpart{Требования к программе}
\begin{itemize}
\item код, полученный от модели, считается недоверенным: перед запуском проверяется по синтаксическому
дереву (модуль \code{ast}: разрешены только импорты \code{numpy} и \code{math}, запрещены \code{open},
\code{exec}, \code{eval}, \code{\_\_import\_\_}, модули \code{os}, \code{sys}, \code{subprocess} и атрибуты
вида \code{\_\_имя\_\_}) и выполняется в отдельном процессе с тайм-аутом; упавшая, зависшая или вернувшая
нечисловую оценку программа отбрасывается;
\item языковая модель подключается функцией \code{llm(prompt) -> (текст, стоимость)}: Claude через
Anthropic SDK или Claude Code CLI (приложение~\ref{app:listings}), модели других провайдеров или локальная
открытая модель~--- по выбору студента;
\item ключ доступа к API задаётся только переменной окружения; он не записывается в код, ноутбук,
отчёт и репозиторий;
\item ведётся журнал вызовов: подсказка, ответ, оценка, стоимость; поиск останавливается по бюджету;
\item отображаются: кривая <<лучшая оценка в базе от числа вызовов>> с уровнями First Fit, Best Fit и
параметрической эвристики; таблица оценок на обучающих и тестовых наборах; код лучшей эвристики.
\end{itemize}

\labpart{Порядок выполнения работы}
\begin{steps}
\item Реализовать генератор наборов, онлайн-упаковку с функцией приоритета и оценку $E$; сверить
First Fit и Best Fit с контрольными значениями таблицы~\ref{tab:l9-var}, объяснить результат Worst Fit.
\item Подобрать параметры $(a, t, b)$ CMA-ES (240 оценок, $\lambda = 8$) и объяснить найденные значения.
\item Реализовать проверку и изолированный запуск кода; убедиться, что программы с \code{import os},
\code{open(...)} и бесконечным циклом отвергаются.
\item Выполнить мини-FunSearch (60 вызовов) и построить кривую поиска; записать число корректных
программ, время и стоимость.
\item Оценить пять лучших эвристик на тестовых наборах и на варианте переноса, выполнить абляцию лучшей
эвристики и описать словами, какие правила она содержит.
\end{steps}

\labpart{Пример выполнения}
Вариант~3: распределение Вейбулла, $C = 100$~--- та же постановка, что в статье FunSearch~\cite{romera2024}.
Модель \code{claude-opus-5-5} (уровень рассуждений \code{medium}) вызывалась через Claude Code CLI в
неинтерактивном режиме без инструментов. Все 60 ответов дали корректный код; поиск занял 5~мин и стоил
2,8 доллара США. Для сравнения тот же поиск выполнен для смеси $0,5\,N(30; 5) + 0,5\,N(65; 5)$ при
$C = 100$: 58 корректных программ из 60, 5,5~мин, 3,1 доллара (рисунок~\ref{fig:l9-progress},
таблица~\ref{tab:l9-transfer}).

\begin{table}[!htb]
\caption{Вариант~3: превышение нижней границы $L_1$ на обучающих ($5 \times 1000$) и тестовых ($5 \times 5000$) наборах, \%}
\label{tab:l9-res}
\centering\small
\begin{tabular}{lcc}
\toprule
Эвристика & Обучение & Тест \\
\midrule
Worst Fit & 13,23 & 13,11 \\
First Fit & 5,25 & 4,36 \\
Best Fit & 4,76 & 3,97 \\
параметрическая, CMA-ES: $a = 45,2$, $t = 24,8$, $b = 8,7$ & 3,17 & 2,29 \\
мини-FunSearch, вызов 25 (эвристика из~\cite{romera2024}) & 3,42 & 0,67 \\
мини-FunSearch, лучшая на обучении (вызов 51) & 2,03 & 0,74 \\
\bottomrule
\end{tabular}
\end{table}

\begin{figure}[!htb]
\centering
\includegraphics[width=\textwidth]{\figdir/lab9_progress.png}
\caption{Лучшая оценка в базе программ от числа вызовов языковой модели (обучающие наборы, $C = 100$):
слева~--- распределение Вейбулла (вариант~3), справа~--- смесь двух нормальных распределений}
\label{fig:l9-progress}
\end{figure}

\emph{Часть А.} Worst Fit оставляет в каждом контейнере большой остаток и втрое хуже остальных правил.
CMA-ES за 240 оценок (9~с) нашла штраф $a = 45$ за остаток меньше $t = 25$: такой остаток почти
никогда не заполнится, поскольку типичный предмет имеет размер около 40. Одна понятная гипотеза и
три параметра улучшают Best Fit в 1,7 раза на тесте.

\emph{Часть Б.} Кривая поиска ступенчатая (рисунок~\ref{fig:l9-progress}): большинство предложений модели
не лучше текущего рекорда, улучшения случаются редко и скачками. На 25-м вызове появилась
<<разностная>> квадратичная оценка из эвристики, опубликованной в статье FunSearch: оценка на обучающих
наборах совпала с опубликованной эвристикой до третьего знака, а её формула дословно вошла в ядро всех
последующих версий. В нашем окружении эта эвристика даёт 0,67\,\% на тесте (в статье~--- 0,68\,\%).
По-видимому, модель знает её из обучающих данных. Такая \emph{утечка} типична для LLM: результат поиска нужно
сверять с литературой, прежде чем называть его новым. Дальнейшие версии добавили к ядру правила,
перечисленные в таблице~\ref{tab:l9-ablation}; почти точным считается остаток не больше 1\,\% ёмкости.

\emph{Часть В.} На тесте лучшая на обучении эвристика уступает своему ядру (0,74 против 0,67\,\%).
Абляция (таблица~\ref{tab:l9-ablation}) объясняет это: правило почти точного заполнения улучшает
оценку на наборах по 1000 предметов почти вдвое, а на наборах по 5000 предметов немного ухудшает;
штраф за <<мёртвый>> остаток не меняет упаковку совсем~--- модель добавила декоративный код.
Поиск подогнал эвристику под условия обучения, как модель машинного обучения подгоняется под обучающую
выборку; поэтому обучающие наборы должны соответствовать условиям применения, а итоговую версию
следует выбирать по отложенным наборам.

\begin{table}[!htb]
\caption{Абляция лучшей эвристики варианта~3: $E$, \%}
\label{tab:l9-ablation}
\centering\small
\begin{tabular}{lcc}
\toprule
Эвристика & Обучение & Тест \\
\midrule
лучшая (все правила) & 2,03 & 0,74 \\
без штрафа за <<мёртвый>> остаток & 2,03 & 0,74 \\
без бонуса за остаток, равный предмету & 2,18 & 0,76 \\
без правила почти точного заполнения & 4,06 & 0,68 \\
без правила точного заполнения & 2,33 & 0,99 \\
только ядро (эвристика из~\cite{romera2024}) & 3,42 & 0,67 \\
\bottomrule
\end{tabular}
\end{table}

Перенос на другие данные (таблица~\ref{tab:l9-transfer}) показывает, что найденные эвристики
специализированы. На ёмкость $C = 150$ при том же распределении выигрыш переносится (1,06 против
2,22\,\% у Best Fit), а на смеси эвристики, выведенные для распределения Вейбулла, в 3,3--3,6 раза
хуже Best Fit. На первом тестовом наборе смеси эвристика из~\cite{romera2024} оставляет в одиночестве
956 из 2531 крупного предмета (Best Fit~--- 104), а мелкие укладывает в отдельные контейнеры примерно
по два.

На смеси поиск почти ничего не дал: одно улучшение за 60 вызовов, на тесте 6,27 против 6,28\,\% у
Best Fit. Причина~--- в нижней границе. Предметы крупнее $C/2$ не помещаются вдвоём, поэтому
контейнеров не меньше, чем таких предметов, и эта граница здесь в среднем на 5,8\,\% выше $L_1$. Best Fit
превышает её всего на 0,46\,\%: улучшать почти нечего. Прежде чем тратить бюджет на поиск, полезно
оценить запас по более точной нижней границе. Параметры CMA-ES на смеси улучшили обучающую оценку
(7,19 против 7,48\,\%), но не тестовую (6,29 против 6,28\,\%)~--- та же подгонка под обучающие наборы.

\begin{table}[!htb]
\caption{Перенос эвристик: $E$ на тестовых наборах ($5 \times 5000$), \%}
\label{tab:l9-transfer}
\centering\small
\begin{tabular}{lccc}
\toprule
Эвристика (где настроена) & \makecell{Вейбулла,\\ $C = 100$} & \makecell{Вейбулла,\\ $C = 150$} & \makecell{смесь,\\ $C = 100$} \\
\midrule
First Fit & 4,36 & 2,40 & 6,45 \\
Best Fit & 3,97 & 2,22 & 6,28 \\
параметрическая + CMA-ES (Вейбулла, 100) & 2,29 & 2,04 & 6,36 \\
эвристика из~\cite{romera2024} & 0,67 & 0,67 & 22,33 \\
мини-FunSearch (Вейбулла, 100) & 0,74 & 1,06 & 21,02 \\
мини-FunSearch (смесь, 100) & 3,97 & 2,25 & 6,27 \\
\bottomrule
\end{tabular}
\end{table}

Итог: поиск с языковой моделью за несколько долларов нашёл эвристику, в 5 раз лучшую Best Fit на
своём распределении (но не на другом), и частично воспроизвёл известный результат. Параметрическая
эвристика с CMA-ES хуже, зато бесплатна, понятна и настраивается за секунды~--- с неё стоит начинать.

\labpart{Контрольные вопросы}
\begin{questions}
\item Чем эволюция программ с языковой моделью отличается от генетического программирования? Какие
операторы ГА заменяет модель?
\item Опишите цикл FunSearch. Зачем в подсказку помещают несколько версий, упорядоченных по возрастанию
качества, и зачем база программ делится на острова?
\item Каким требованиям должен удовлетворять оценщик? Почему код модели выполняется в отдельном процессе
с тайм-аутом и почему статической проверки кода недостаточно в промышленной системе?
\item Сравните правила First Fit, Best Fit и Worst Fit. Запишите нижнюю границу $L_1$; почему для
смеси крупных и мелких предметов она недостижима?
\item Когда достаточно параметрической эвристики с настройкой CMA-ES, а когда нужен поиск в пространстве
программ?
\item Почему улучшение на обучающих наборах может не переноситься на тестовые? Как выбирать итоговую
эвристику?
\item Что такое утечка данных при поиске с языковой моделью и как её обнаружить? Как учитывать
стоимость вызовов модели при сравнении методов?
\end{questions}
